\chapter{Apprentissage automatique et ingénieur de données}\label{ch:in:machine-learning-and-data-engineer}

Pour l'exercice 2021, les employeurs financiers retenus ont déposé 136 dossiers de conditions de travail (LCA) pour des emplois intitulés en apprentissage automatique,
data science ou ingénierie des données~; pour l'exercice 2025, ils en ont déposé 1\,051, soit près de huit fois plus. Sur ces mêmes quatre années
leurs dépôts pour les ingénieurs logiciels ont doublé et ceux pour les \omterm{def:in:quant-researcher:def}{chercheurs quantitatifs} n'ont pas bougé. La profession que la plupart des dépôts de 2025
désignent, \omterm{def:in:machine-learning-and-data-engineer:roles}{data scientist}, n'existait pas dans la classification américaine de 2010~; elle a été créée lors de la révision de 2018.
Ce chapitre décrit les emplois en apprentissage automatique et en données dans le trading et la finance, mesure à quelle vitesse leurs dépôts
ont progressé et avec quelle incertitude, et compare leur rémunération à celle des ingénieurs logiciels qui leur sont voisins.

\begin{center}\footnotesize
\begin{tabular}{@{}p{2.4cm}p{3.3cm}p{3.3cm}p{3.3cm}@{}}
\toprule
\multicolumn{4}{@{}l}{\textbf{Fiches de rôle~: apprentissage automatique et données}}\\
\midrule
 & \omterm{def:in:machine-learning-and-data-engineer:roles}{data scientist} & ingénieur en apprentissage automatique & \omterm{def:in:machine-learning-and-data-engineer:roles}{ingénieur de données}\\
\midrule
travaille sur & modèles tirés des données, pour une question de recherche ou métier & les systèmes qui entraînent, servent et surveillent les modèles & les chaînes de traitement et les bases où résident les données\\
se situe dans & groupes de recherche ou lignes métier & une équipe plateforme ou un groupe de recherche & une équipe données ou plateforme\\
codes dans les dépôts & 15-2051 & 15-2051, 15-1252 & 15-1252, 15-1243\\
traité dans & livre~12, ch.~1~; livre~7, ch.~12 & livre~12, ch.~23--28 & livre~15, ch.~2--4~; livre~12, ch.~24\\
\bottomrule
\end{tabular}
\end{center}

\begin{definition}[Data scientist, ingénieur de données]\label{def:in:machine-learning-and-data-engineer:roles}
Un \emph{data scientist}\index{data scientist} est un analyste qui construit des modèles statistiques et d'apprentissage automatique à partir de données
pour répondre à une question ou guider une décision, et rend compte de ce que montrent les données. Un \emph{ingénieur de données}\index{ingénieur de données} (data engineer)
est un ingénieur qui conçoit et exploite les chaînes de traitement, le stockage et les interfaces par lesquels les données sont acquises, nettoyées,
versionnées et mises à disposition des personnes et des systèmes qui les utilisent.
\end{definition}

La définition professionnelle du premier est large~: les \omterm{def:in:machine-learning-and-data-engineer:roles}{data scientists} «~Develop and implement a set of techniques or
analytics applications to transform raw data into meaningful information~» et «~Apply data mining, data modeling,
natural language processing, and machine learning~». La classification ne prévoit aucune profession pour le second~: O*NET
mentionne «~Data Engineer~» parmi les intitulés associés aux architectes de bases de données (15-1243), et les employeurs déclarent la plupart des \omterm{def:in:machine-learning-and-data-engineer:roles}{ingénieurs de données} comme
développeurs logiciels. L'ingénieur en apprentissage automatique (ML engineer), qui construit les systèmes qui entraînent, servent et surveillent les modèles, est
traité dans le livre~12 (chapitre~28).

\section{Rôles en apprentissage automatique intégrés à la recherche}

Dans le groupe de recherche d'une société de trading, l'apprentissage automatique est une méthode, non un département. Le chercheur du chapitre~17 qui ajuste
un modèle à gradient boosting aux variables du carnet d'ordres fait de l'apprentissage automatique~; l'intitulé du contrat peut dire
\omterm{def:in:quant-researcher:def}{chercheur quantitatif} (quant researcher) ou \omterm{def:in:machine-learning-and-data-engineer:roles}{data scientist}, et les dépôts suivent l'intitulé (chapitre~19). Ce qui distingue le
rôle intégré à la recherche est le problème~: un signal au faible rapport signal sur bruit, peu d'observations indépendantes,
des données non stationnaires et un coût de l'erreur qui se paie en argent. Le livre~12 porte précisément sur cette différence
(chapitre~1), et ses outils (validation qui respecte le temps, étiquettes construites à partir de prix futurs, modèles qui résistent au bruit)
sont le quotidien du rôle.

Un spécialiste de l'apprentissage automatique intégré à la recherche est jugé comme l'est un chercheur~: sur ses résultats hors échantillon (out-of-sample), et sur le fait que le
modèle apporte par rapport à ce que font déjà des modèles plus simples. L'habitude du livre~12 de confronter chaque modèle à un modèle linéaire régularisé
de référence (chapitre~4) est la discipline du rôle. La \omterm{def:in:quant-researcher:feedback}{vitesse du retour d'information} du chapitre~17 s'applique telle quelle~: un modèle rééquilibré chaque mois
accumule des preuves aussi lentement que n'importe quel signal mensuel.

\section{Rôles de plateforme en apprentissage automatique}

Une société qui fait tourner de nombreux modèles a besoin d'une plateforme d'apprentissage automatique (livre~12, chapitre~28)~: entrepôts de données et feature stores, infrastructure d'entraînement,
suivi des expériences et registre de modèles, compilation et mise en service, surveillance et réentraînement. L'ingénieur en
apprentissage automatique le construit et l'exploite, et la spécification de transmission (handoff) du livre~12 est le contrat entre les
personnes qui construisent les modèles et celles qui les exploitent. La plus grande part de la croissance des dépôts (ci-dessous) provient des banques, et les
dépôts ne précisent pas à laquelle des activités de la banque chaque emploi se rattache.

Le rôle de plateforme est jugé comme celui de l'\omterm{def:in:quant-developer:research}{ingénieur de recherche} du chapitre~19~: sur la rapidité et la sûreté avec lesquelles les modèles atteignent
la production et sur la rareté de leurs défaillances à ce stade. Il est plus proche du génie logiciel que de la recherche, et sa rémunération et ses
intitulés suivent davantage ceux du génie logiciel que ceux de la recherche.

\section{L'ingénieur de données}

Une société de trading est d'abord une consommatrice de données~: données de marché par plateforme et par jour, données de référence,
opérations sur titres, données alternatives (livre~7, chapitre~12), ses propres ordres et exécutions. L'\omterm{def:in:machine-learning-and-data-engineer:roles}{ingénieur de données} gère le trajet depuis
le fichier du fournisseur ou de la plateforme jusqu'à la requête du chercheur~: capture et stockage des ticks (livre~15, chapitre~2), formats en colonnes
et bases de ticks (tick stores) (chapitres~3 et~4), versions point-in-time qui permettent à un backtest de ne voir que ce qui était connu
au moment considéré, des feature stores qui servent les mêmes valeurs à l'entraînement et à la production (livre~12, chapitre~24), et les contrôles qui
détectent une modification silencieuse du fournisseur avant qu'un modèle ne s'y entraîne.

Ce travail est peu flatteur mais indispensable~: un résultat de recherche ne vaut que par les données sur lesquelles il repose, et une journée manquante, un
horodatage erroné ou un univers biaisé par la survie peuvent produire un backtest faux dans chaque détail. Les \omterm{def:in:machine-learning-and-data-engineer:roles}{ingénieurs de données} sont
jugés sur l'exhaustivité, l'exactitude et la ponctualité de ce qu'ils servent, et sur le coût de leur stockage et de leur transfert.

\section{Croissance et rémunération}

\begin{dated}{2025-09}{Rôles d'apprentissage automatique et de données dans les dépôts}\label{dat:in:machine-learning-and-data-engineer:filings}
Dossiers LCA des employeurs financiers retenus, exercices 2021 et 2025~: intitulés en apprentissage automatique et en données
136 et 1\,051 (dont 920 pour les banques en 2025)~; ingénieurs logiciels 2\,618 et 5\,727~; \omterm{def:in:quant-researcher:def}{chercheurs quantitatifs} 1\,287 et 1\,297~;
toutes familles 4\,422 et 8\,808. Pour l'exercice 2025, \omterm{def:in:machine-learning-and-data-engineer:roles}{data scientists} (15-2051) contre les développeurs logiciels (15-1252) chez les mêmes
employeurs~: médiane du \omterm{def:in:how-pay-works:base}{salaire de base} proposé de 138\,000~\$ contre 155\,000~\$~; par \omterm{def:in:pay-levels-by-role-firm-type-and-seniority:soc}{niveau de salaire} $-\$20\,500$ (I), $-\$18\,300$ (II),
$-\$19\,000$ (III), $+\$2\,220$ (IV). \omterm{def:in:how-pay-works:base}{Salaire de base} seulement.
\end{dated}

\begin{method}[Croissance des nombres de dépôts]\label{met:in:machine-learning-and-data-engineer:trend}
Soit $n_t$ le nombre de dépôts de l'année $t$. On les modélise comme une loi de Poisson de moyenne $\lambda_0e^{g t}$~; l'ajustement par moindres carrés pondérés
de $\log n_t$ en $t$, avec pour poids $n_t$, estime le taux de croissance continu $g$, dont l'erreur type vaut
$1/\sqrt{\sum_t n_t(t-\bar t)^2}$. Avec deux années $t_0,t_1$, on a $g=\log(n_1/n_0)/(t_1-t_0)$ et
une erreur type de $\sqrt{1/n_0+1/n_1}/(t_1-t_0)$, et le taux annuel est $e^g-1$. Chaque nombre admet un intervalle de Poisson exact.
\end{method}

Les dépôts en apprentissage automatique et en données ont progressé de 66.7~\% par an (intervalle à 95~\% de 59.4~\% à 74.3~\%), ceux des ingénieurs logiciels à
21.6~\% (20.2~\% à 23.0~\%)~; celui des \omterm{def:in:quant-researcher:def}{chercheurs quantitatifs} n'a pas progressé (0.2~\%, $-1.7\%$ à 2.1~\%) et celui des traders a
légèrement baissé (\cref{fig:in:machine-learning-and-data-engineer:growth}). La part de la famille dans l'ensemble des dépôts est passée de 3.1~\% à
11.9~\%. L'intervalle de Poisson est le plus petit intervalle honnête~: les dépôts arrivent par lots de quelques employeurs, et neuf sur dix des
dépôts de 2025 en apprentissage automatique viennent de banques, si bien que l'incertitude réelle sur les recrutements du secteur est plus grande que
ce que ne le laisse croire l'intervalle.

\begin{omfigure}
\begin{tikzpicture}
\begin{axis}[ybar,area legend,width=10.5cm,height=5.4cm,xmin=0.5,xmax=6.5,ymode=log,ymin=10,ymax=10000,log origin=infty,bar width=11pt,
  xtick={1,2,3,4,5,6},xticklabels={{ML et données},{ingénieur logiciel},{développeur quant},{risque},{chercheur quantitatif},{trader}},
  xticklabel style={font=\scriptsize,text width=1.6cm,align=center},ylabel={\small dépôts (échelle log)},
  yticklabel style={font=\footnotesize},grid=major,grid style={gray!20},table/col sep=comma,
  legend style={legend cell align=left,font=\scriptsize,at={(0.5,-0.42)},anchor=north,legend columns=2,/tikz/every even column/.append style={column sep=8pt}}]
\addplot[fill=black!30,draw=none,error bars/.cd,y dir=both,y explicit]
  table[x=pos,y=n2021,y error plus expr=\thisrow{hi2021}-\thisrow{n2021},y error minus expr=\thisrow{n2021}-\thisrow{lo2021}]
  {figdata/industry/21-machine-learning-and-data-engineer/counts.csv};
\addplot[fill=omDef,draw=none,error bars/.cd,y dir=both,y explicit]
  table[x=pos,y=n2025,y error plus expr=\thisrow{hi2025}-\thisrow{n2025},y error minus expr=\thisrow{n2025}-\thisrow{lo2025}]
  {figdata/industry/21-machine-learning-and-data-engineer/counts.csv};
\legend{exercice 2021,exercice 2025}
\end{axis}
\end{tikzpicture}

\omcaption{Effectifs des demandes de certification du travail par famille de métiers chez les employeurs financiers sourcés, exercices 2021 et 2025, avec
des intervalles de Poisson exacts à 95~\% (à peine visibles à ces effectifs). Données~: \texttt{data/industry/lca\_ranges.csv}, via
\texttt{in\_mldata.counts}.}\label{fig:in:machine-learning-and-data-engineer:counts}
\end{omfigure}

\begin{omfigure}
\begin{tikzpicture}
\begin{axis}[width=10.5cm,height=5cm,xmin=-20,xmax=90,ymin=0.4,ymax=6.6,ytick=data,
  yticklabels={{trader},{chercheur quantitatif},{risque},{développeur quant},{ingénieur logiciel},{ML et données}},
  yticklabel style={font=\footnotesize},xticklabel style={font=\footnotesize},
  xlabel={\small taux de croissance annuel des dépôts, exercice 2021 à exercice 2025, \%},grid=major,grid style={gray!20},table/col sep=comma]
\addplot[only marks,mark=*,omDef,error bars/.cd,x dir=both,x explicit]
  table[x=annual,y=pos,x error plus expr=\thisrow{hi}-\thisrow{annual},x error minus expr=\thisrow{annual}-\thisrow{lo}]
  {figdata/industry/21-machine-learning-and-data-engineer/growth.csv};
\draw[black!50] (axis cs:0,0.4) -- (axis cs:0,6.6);
\end{axis}
\end{tikzpicture}

\omcaption{Taux de croissance annuel des dépôts par famille de métiers entre l'exercice 2021 et l'exercice 2025, avec des intervalles à 95~\% selon le
modèle de Poisson. Données~: \texttt{in\_mldata.trends}, issues de \texttt{firm.roles.filing\_trend}.}\label{fig:in:machine-learning-and-data-engineer:growth}
\end{omfigure}

La comparaison des salaires est moins flatteuse que la croissance. Chez les mêmes employeurs, les dépôts classés comme \omterm{def:in:machine-learning-and-data-engineer:roles}{data scientists} affichent un
\omterm{def:in:how-pay-works:base}{salaire de base} proposé dont la médiane est inférieure de 17\,000~\$ à celle des développeurs logiciels (intervalle à 95~\%~: $-\$20\,000$ à $-\$15\,000$), et
d'environ 19\,000~\$ de moins à chacun des niveaux I à III~; seul le niveau IV rapproche les deux (\cref{fig:in:machine-learning-and-data-engineer:gap}).
L'essentiel de la croissance se trouve dans les banques, dont les dépôts ne précisent pas à quelle activité sert le poste~; les rares
dépôts d'apprentissage automatique émanant de sociétés de trading affichent des médianes plus élevées (175\,000~\$ chez les teneurs de marché, 193\,000~\$ sur les plateformes et
192\,912~\$ chez les fonds systématiques, chacune à partir de dix à douze dépôts), en cohérence avec les chercheurs et ingénieurs de ces sociétés.

\begin{omfigure}
\begin{tikzpicture}
\begin{axis}[width=10.5cm,height=5cm,xmin=0.5,xmax=5.5,ymin=-30,ymax=10,xtick={1,2,3,4,5},
  xticklabels={tous niveaux,I,II,III,IV},xlabel={\small niveau de salaire},ylabel={\small écart, en milliers de \$},
  tick label style={font=\footnotesize},grid=major,grid style={gray!20},table/col sep=comma]
\addplot[only marks,mark=*,omDef,thick,error bars/.cd,y dir=both,y explicit]
  table[x=pos,y=diff,y error plus expr=\thisrow{hi}-\thisrow{diff},y error minus expr=\thisrow{diff}-\thisrow{lo}]
  {figdata/industry/21-machine-learning-and-data-engineer/gap.csv};
\draw[black!50] (axis cs:0.5,0) -- (axis cs:5.5,0);
\end{axis}
\end{tikzpicture}

\omcaption{\omterm{def:in:how-pay-works:base}{Salaire de base} médian proposé des dépôts de \omterm{def:in:machine-learning-and-data-engineer:roles}{data scientists} moins celui des développeurs logiciels chez les employeurs financiers sourcés,
exercice 2025, pour l'ensemble et par \omterm{def:in:pay-levels-by-role-firm-type-and-seniority:soc}{niveau de salaire}, avec des intervalles bootstrap à 95~\%. Données~: \texttt{data/industry/lca\_ds\_gap.csv},
via \texttt{in\_mldata.gaps}.}\label{fig:in:machine-learning-and-data-engineer:gap}
\end{omfigure}

\begin{dated}{2025-05}{Où travaillent les data scientists~: l'enquête}\label{dat:in:machine-learning-and-data-engineer:survey}
Enquête professionnelle, mai 2025, \omterm{def:in:machine-learning-and-data-engineer:roles}{data scientists} (15-2051)~: finance et assurance 46\,730 salariés, médiane 124\,770~\$~; information 32\,410, médiane 141\,440~\$~; services professionnels, scientifiques et techniques 69\,730, médiane 126\,730~\$.
information 32\,410, médiane 141\,440~\$~; services professionnels, scientifiques et techniques 69\,730, médiane 126\,730~\$.
Au sein de la finance~: intermédiation de crédit 12\,450 (médiane 130\,300~\$), assureurs 15\,090 (107\,680~\$), valeurs mobilières
7\,210 (134\,510~\$). Éditeurs de logiciels 10\,950 (156\,220~\$).
\end{dated}

La finance et l'assurance emploient davantage de \omterm{def:in:machine-learning-and-data-engineer:roles}{data scientists} que le secteur de l'information, mais à une médiane plus basse, et
l'activité des valeurs mobilières n'en représente qu'une faible part (\cref{fig:in:machine-learning-and-data-engineer:survey}). Ce constat
est confirmé par les dépôts~: la science des données dans la finance est surtout un métier de banque et d'assurance, et dans le trading elle est
surtout l'un des outils de la recherche.

\begin{omfigure}
\begin{tikzpicture}
\begin{axis}[xbar,width=10.5cm,height=5.4cm,xmin=0,xmax=85,ymin=0.4,ymax=7.6,bar width=9pt,ytick=data,
  yticklabels={{éditeurs de logiciels},{valeurs mobilières},{banques (intermédiation de crédit)},{assureurs},{information},{finance et assurance},{services professionnels et techniques}},
  yticklabel style={font=\scriptsize},xticklabel style={font=\footnotesize},
  xlabel={\small data scientists employés, en milliers},grid=major,grid style={gray!20},table/col sep=comma,
  nodes near coords,nodes near coords style={font=\scriptsize,anchor=west,xshift=1pt,yshift=0pt},
  point meta=explicit symbolic]
\addplot[fill=omDef,draw=none] table[x=employment,y=pos,meta=median]
  {figdata/industry/21-machine-learning-and-data-engineer/survey.csv};
\end{axis}
\end{tikzpicture}

\omcaption{\omterm{def:in:machine-learning-and-data-engineer:roles}{Data scientists} employés par secteur et par activité, enquête professionnelle de mai 2025~; les étiquettes donnent le salaire annuel
médian en milliers de \$. Les trois premières barres sont des secteurs~; les quatre en dessous sont des activités à l'intérieur de ceux-ci (banques, assureurs
et valeurs mobilières au sein de la finance et de l'assurance~; éditeurs de logiciels au sein de l'information). Données~:
\texttt{data/industry/oews\_roles.csv}, via \texttt{in\_mldata.survey}.}\label{fig:in:machine-learning-and-data-engineer:survey}
\end{omfigure}

\section{Tutoriel~: le titre qui croît le plus vite}\label{tut:in:machine-learning-and-data-engineer:tut}

\begin{tutorial}
\textbf{Objectif.} Mesurer la croissance des dépôts d'apprentissage automatique et de données avec leur incertitude, et mettre leurs salaires en regard
de ceux des ingénieurs logiciels.
\textbf{État final~:} les \cref{fig:in:machine-learning-and-data-engineer:counts,fig:in:machine-learning-and-data-engineer:growth,fig:in:machine-learning-and-data-engineer:gap}.
\begin{tutsteps}
\item \textbf{Effectifs.} Sommer le fichier \texttt{lca\_ranges.csv} du chapitre~14 par type d'employeur pour chaque famille de métiers et chaque année~;
\texttt{firm.roles.poisson\_interval(n)} donne l'intervalle exact de chaque effectif.
\item \textbf{Croissance.} \texttt{filing\_trend(years, counts)} ajuste le modèle de Poisson log-linéaire
(\cref{lst:in:machine-learning-and-data-engineer:trend}).
\omcode{code/firm/roles/firm_roles.py}{320}{331}{Croissance des effectifs de dépôts avec son erreur type.}\label{lst:in:machine-learning-and-data-engineer:trend}
\item \textbf{Salaires.} \texttt{in\_lca\_ds\_derive.py} compare, niveau par niveau, les dépôts de \omterm{def:in:machine-learning-and-data-engineer:roles}{data scientists} et de développeurs logiciels
issus du cache du chapitre~14, avec \texttt{median\_gap} (chapitre~19).
\item \textbf{Enquête.} \texttt{oews\_roles.csv} pour 15-2051 par secteur et par activité.
\end{tutsteps}
\end{tutorial}

Pour la famille de l'apprentissage automatique~: $g=\log(1\,051/136)/4=0.511$ avec une erreur type de 0.023~; soit 66.7~\% par an. L'intervalle de Poisson de son effectif
2025 va de 988 à 1\,117.

\textbf{Pour aller plus loin.} Lire les autres exercices fiscaux et ajuster la tendance sur l'ensemble~; séparer les banques des sociétés de trading~;
remplacer le modèle de Poisson par un modèle qui tient compte des employeurs qui déposent par lots (une binomiale négative, ou un bootstrap
sur les employeurs).

\section{Construction~: cartes d'apprentissage automatique et tendances des dépôts}\label{bld:in:machine-learning-and-data-engineer:roles}

\begin{build}
\bfield{Rôle} Ajouter les rôles d'apprentissage automatique et de données au registre, et mesurer la croissance d'un effectif avec
son incertitude.
\bfield{Interface} \texttt{firm.roles}~: les cartes \texttt{data scientist}, \texttt{data engineer},
\texttt{machine-learning engineer}~; \texttt{poisson\_interval(n, level) -> (lo, hi)}~;
\texttt{filing\_trend(years, counts, z) -> dict(rate, se, annual, lo, hi)}.
\bfield{Règles} Les effectifs suivent un modèle de Poisson sauf indication contraire de l'appelant~; l'intervalle est exact (Garwood)~; la croissance est
continue dans l'ajustement et rapportée comme un taux de croissance annuel composé.
\bfield{Tests d'acceptation} \texttt{code/firm/roles/tests/}~: deux points donnent la forme close~; des effectifs doublant chaque année donnent
$\log 2$ exactement~; l'intervalle pour dix va de 4.80 à 18.39~; zéro donne une borne inférieure nulle.
\bfield{Pour aller plus loin} Surdispersion~; une tendance par type d'employeur avec un test de croissance égale~; prévisions avec intervalles.
\end{build}

\begin{omsources}
\item Bureau of Labor Statistics, SOC 2010 to 2018 crosswalk~; enquête sur les professions, mai 2025.
\item O*NET OnLine, professions 15-2051 et 15-1243 (CC BY 4.0).
\item Tableaux du chapitre~14, issus des fichiers LCA du Department of Labor.
\item Garwood, F. (1936), Fiducial limits for the Poisson distribution, \emph{Biometrika} 28(3/4), 437--442.
\end{omsources}

\section{Exercices}

\begin{exercise}[$\star$]\label{exo:in:machine-learning-and-data-engineer:1}
Déterminer de quel facteur les dépôts d'apprentissage automatique et de données ont augmenté entre l'exercice~2021 et 2025, et
exprimer ce résultat sous forme de taux annuel.
\end{exercise}

\begin{exercise}[$\star$]\label{exo:in:machine-learning-and-data-engineer:2}
Déterminer la part que représente cette famille dans l'ensemble des dépôts, pour chaque année.
\end{exercise}

\begin{exercise}[$\star$]\label{exo:in:machine-learning-and-data-engineer:3}
À partir de l'encadré de l'enquête, déterminer quel secteur emploie le plus de \omterm{def:in:machine-learning-and-data-engineer:roles}{data scientists}, la finance et l'assurance ou le secteur de l'information, et
lequel verse le salaire médian le plus élevé.
\end{exercise}

\begin{exercise}[$\star\star$]\label{exo:in:machine-learning-and-data-engineer:4}
Calculer l'erreur type du taux de croissance pour les ingénieurs logiciel et expliquer pourquoi elle est
plus faible que celle de la famille d'apprentissage automatique.
\end{exercise}

\begin{exercise}[$\star\star$]\label{exo:in:machine-learning-and-data-engineer:5}
Expliquer pourquoi la croissance des dépôts codés 15-2051 ne peut pas être mesurée à partir de l'exercice~2021, alors que celle de la famille de titres peut l'être.
\end{exercise}

\begin{exercise}[$\star\star$]\label{exo:in:machine-learning-and-data-engineer:6}
Un diplômé compare une offre de \omterm{def:in:machine-learning-and-data-engineer:roles}{data scientist} dans une banque à une offre de \omterm{def:in:quant-researcher:def}{chercheur quantitatif} (quant researcher) dans un fonds systématique.
Dire ce que disent les éléments du chapitre, et ce qu'ils ne disent pas.
\end{exercise}

\begin{exercise}[$\star\star\star$]\label{exo:in:machine-learning-and-data-engineer:7}
\emph{Programmation.} Supposer que les dépôts d'apprentissage automatique des deux années arrivent par lots de dix, de sorte que
chaque effectif se comporte comme dix fois un comptage de Poisson. Recalculer l'erreur type et l'intervalle du taux de croissance.
\end{exercise}

\begin{exercise}[$\star\star\star$]\label{exo:in:machine-learning-and-data-engineer:8}
\emph{Trouver la faille.} «~Les dépôts d'apprentissage automatique ont progressé de 67~\% par an alors que ceux des chercheurs n'ont pas progressé, donc les sociétés de trading
remplacent les \omterm{def:in:quant-researcher:def}{chercheurs quantitatifs} par des \omterm{def:in:machine-learning-and-data-engineer:roles}{data scientists}.~»
\end{exercise}

\section{Problème~: L'intitulé à la croissance la plus rapide}

\begin{problem}[{Problème du week-end --- l'intitulé à la croissance la plus rapide}]\label{pb:in:machine-learning-and-data-engineer:1}
Un étudiant qui hésite entre un master de statistique avec un parcours en apprentissage automatique et un master de mathématiques
financières se demande dans quel sens évolue le recrutement du secteur.

\medskip\noindent\textbf{Partie~I --- Les rôles.}
\begin{enumerate}
\item Définir le \omterm{def:in:machine-learning-and-data-engineer:roles}{data scientist} et l'\omterm{def:in:machine-learning-and-data-engineer:roles}{ingénieur de données} (data engineer).
\item Qu'est-ce qui distingue un rôle d'apprentissage automatique intégré à la recherche~?
\item Que construit l'ingénieur en apprentissage automatique (ML engineer), et qu'est-ce que la spécification de transmission (handoff)~?
\item Que prend en charge un \omterm{def:in:machine-learning-and-data-engineer:roles}{ingénieur de données} dans une société de trading~?
\item Comment la classification traite-t-elle les \omterm{def:in:machine-learning-and-data-engineer:roles}{ingénieurs de données}~?
\end{enumerate}

\medskip\noindent\textbf{Partie~II --- Les effectifs.}
\begin{enumerate}[resume]
\item Donner les effectifs de la famille pour les exercices 2021 et 2025.
\item Énoncer le modèle de croissance et son erreur type.
\item Donner les taux de croissance annuels, avec leurs intervalles, pour la famille d'apprentissage automatique, pour les ingénieurs logiciel et pour les chercheurs.
\item Pourquoi l'intervalle de Poisson est-il le plus petit intervalle honnête~?
\item Quel type d'employeur a déposé la majeure partie des dépôts d'apprentissage automatique de 2025~?
\end{enumerate}

\medskip\noindent\textbf{Partie~III --- La rémunération.}
\begin{enumerate}[resume]
\item Donner l'écart de salaire des \omterm{def:in:machine-learning-and-data-engineer:roles}{data scientists} par rapport aux développeurs logiciel, dans l'ensemble et par niveau.
\item Que montrent les rares dépôts d'apprentissage automatique des sociétés de trading~?
\item Où travaillent les \omterm{def:in:machine-learning-and-data-engineer:roles}{data scientists}, selon l'enquête~?
\item Pourquoi le \omterm{def:in:machine-learning-and-data-engineer:roles}{data scientist} d'une banque pourrait-il être moins bien payé que son développeur logiciel, au même niveau~?
\item Que laisse de côté le \omterm{def:in:how-pay-works:base}{salaire de base} ici~?
\end{enumerate}

\medskip\noindent\textbf{Partie~IV --- Le verdict.}
\begin{enumerate}[resume]
\item Énoncer le \emph{résultat nommé}~: le taux de croissance annuel des dépôts d'apprentissage automatique et de données chez les employeurs de la finance,
de l'exercice 2021 à 2025, avec son intervalle, et l'écart du \omterm{def:in:how-pay-works:base}{salaire de base} médian par rapport aux développeurs logiciel, au même \omterm{def:in:pay-levels-by-role-firm-type-and-seniority:soc}{niveau de salaire}.
\item La croissance est-elle liée au trading~?
\item Quel master prépare à quels métiers~?
\item Que doit demander l'étudiant à un employeur au sujet d'un rôle en data science~?
\item En deux phrases, répondre à l'étudiant.
\end{enumerate}
\end{problem}

\section{Questions d'entretien}

\begin{interviewq}[$\star$ \iqroles{ingénieur ML, chercheur}]\label{iq:in:machine-learning-and-data-engineer:1}
Pourquoi la validation croisée k-fold avec mélange est-elle incorrecte pour un modèle de prédiction de rendements quotidiens~?
\end{interviewq}

\begin{interviewq}[$\star$ \iqroles{ingénieur ML}]\label{iq:in:machine-learning-and-data-engineer:2}
La précision en production d'un modèle est inférieure à sa précision en validation. Énumérez cinq causes, dans l'ordre où vous les vérifieriez.
\end{interviewq}

\begin{interviewq}[$\star\star$ \iqroles{ingénieur ML, développeur}]\label{iq:in:machine-learning-and-data-engineer:3}
Comment vous assurez-vous qu'une variable a la même valeur à l'entraînement et en production~?
\end{interviewq}

\begin{interviewq}[$\star\star$ \iqroles{développeur}]\label{iq:in:machine-learning-and-data-engineer:4}
Un fournisseur modifie sans prévenir la définition d'un champ. Comment votre pipeline le remarquerait-il~?
\end{interviewq}

\begin{interviewq}[$\star\star$ \iqroles{ingénieur ML, chercheur}]\label{iq:in:machine-learning-and-data-engineer:5}
Votre modèle à gradient boosting dépasse un modèle linéaire de 0.3~\% de précision. Vaut-il la peine d'être déployé~?
\end{interviewq}

\begin{interviewq}[$\star\star\star$ \iqroles{ingénieur ML}]\label{iq:in:machine-learning-and-data-engineer:6}
Concevez le stockage de dix ans de données de carnet d'ordres que les chercheurs interrogent par symbole et par plage temporelle.
\end{interviewq}
